% ----------------------------------------------------------------
\documentclass[conference]{IEEEtran}
\usepackage{amsmath,amssymb}
\usepackage[T1]{fontenc}
\usepackage[utf8]{inputenc}
\usepackage{lmodern}
\usepackage{microtype}
\usepackage{xcolor}
\usepackage{listings}
\usepackage{enumitem}
\usepackage{tabularx}
\usepackage{graphicx}
\usepackage{titlesec}
\usepackage{xurl}
\usepackage[hidelinks]{hyperref}
% ----------------------------------------------------------------





% ----------------------------------------------------------------

\lstdefinelanguage{SystemVerilog}{
  morekeywords={always,always_ff,always_comb,assert,assign,begin,end,endproperty,
    if,else,property,logic,for},
  sensitive=true,
  morecomment=[l]{//}
}
\lstset{
  language=SystemVerilog,
  basicstyle=\ttfamily\scriptsize,
  columns=fullflexible,
  keepspaces=true,
  frame=single,
  rulecolor=\color{black!55},
  backgroundcolor=\color{black!3},
  framesep=3pt,
  xleftmargin=4pt,
  xrightmargin=4pt,
  breaklines=true,
  showstringspaces=false,
  aboveskip=0.8em,
  belowskip=0.8em
}
\titleformat{\section}
  {\normalfont\large\bfseries}
  {\thesection.}{0.55em}{}
\titleformat{\subsection}
  {\normalfont\normalsize\bfseries}
  {\thesubsection}{0.6em}{}
\titlespacing*{\section}{0pt}{2.2ex plus 0.7ex minus 0.3ex}{1.0ex}
\titlespacing*{\subsection}{0pt}{1.8ex plus 0.5ex minus 0.2ex}{0.7ex}
\setlist[itemize]{
  label=\textbullet,
  leftmargin=1.35em,
  labelsep=0.55em,
  topsep=0.55em,
  itemsep=0.55em,
  parsep=0pt,
  partopsep=0pt
}
\raggedbottom
\setlength{\emergencystretch}{2em}
\setcounter{secnumdepth}{2}
\renewcommand{\thesection}{\arabic{section}}
\renewcommand{\thesubsection}{\thesection.\arabic{subsection}}
\newcommand{\bugsection}[2]{%
  \refstepcounter{section}%
  \setcounter{subsection}{0}%
  \section*{Bug#1: #2}%
}
\urlstyle{same}
\hypersetup{
  hidelinks,
  pdftitle={SystemVerilog RTL Bug Study},
  pdfauthor={Geeth Sathsara Wijesekara}}

% ----------------------------------------------------------------










% ----------------------------------------------------------------

% Title + Cover Page
\title{SystemVerilog RTL Bug Study}
\author{Geeth Sathsara Wijesekara}
\date{}


\begin{document}
\maketitle

% ----------------------------------------------------------------













% ----------------------------------------------------------------

% Introduction

\section*{Introduction}

This report examines five documented RTL design bugs from the open-source
projects Corundum~\cite{corundum}, OpenTitan~\cite{opentitan}, and
CVA6~\cite{cva6}. Each bug is reduced to a minimal reproducer and implemented
in both SystemVerilog and Anvil~\cite{anvil-paper,anvil-github}, then verified
using standard industry approaches and the Anvil compiler where applicable.
The report discusses how these bugs escaped their original verification
environments and whether Anvil's compiler checks could have detected them
earlier.

% ----------------------------------------------------------------




























% ----------------------------------------------------------------
% Bug1
\bugsection{1}{Partial queue-memory initialization}
% -----------------------



% -----------------------
% 1. Instance
\subsection{Instance}

This bug comes from Corundum~\cite{corundum}, an open-source
high-performance FPGA-based network interface card (NIC).

The three affected queue-management modules~\cite{queue-manager-src,
cpl-queue-manager-src,tx-scheduler-src} store per-queue state in inferred
FPGA RAM. During initialization, every RAM entry must be set to zero so that
every legal queue index begins in an empty state.

In commit \texttt{0560f98e}~\cite{corundum-intro}, a split nested loop was added
to support more than 4,000 queues while working around Quartus's loop-iteration
limit. Quartus Prime is Intel's FPGA toolchain~\cite{quartus}. This workaround
introduced the bug discussed below;
it was later corrected in PR~\#96~\cite{corundum-fix} and its fixing
commit~\cite{corundum-fix-commit}.
% -----------------------



% -----------------------
% 2. Mechanism
\subsection{Mechanism}

% - Cause
The above workaround confused the address width with the number of entries. The RAM
contains \texttt{2**QUEUE\_INDEX\_WIDTH} entries, but the new outer loop
stopped at \texttt{QUEUE\_INDEX\_WIDTH}.

\begin{lstlisting}
// Buggy code
for (i = 0; i < QUEUE_INDEX_WIDTH;
     i = i + 2**(QUEUE_INDEX_WIDTH/2))
  for (j = i; j < i + 2**(QUEUE_INDEX_WIDTH/2); j = j + 1)
    queue_ram[j] = 0;
\end{lstlisting}

Therefore, some entries remain uninitialized. The fix changes the outer bound
to \texttt{2**QUEUE\_INDEX\_WIDTH}.

\begin{lstlisting}
// Fixed code
for (i = 0; i < 2**QUEUE_INDEX_WIDTH;
     i = i + 2**(QUEUE_INDEX_WIDTH/2))
  for (j = i; j < i + 2**(QUEUE_INDEX_WIDTH/2); j = j + 1)
    queue_ram[j] = 0;
\end{lstlisting}

% - Minimal Repro
The minimal buggy reproducers in SystemVerilog~\cite{bug1-sv-buggy} and
Anvil~\cite{bug1-anvil-buggy}, together with their fixed SystemVerilog
\cite{bug1-sv-fixed} and Anvil~\cite{bug1-anvil-fixed} versions, are available
in the cited reproduction files.
% -----------------------



% -----------------------
% 3. Why it survived
\subsection{Why It Survived}

% - Why simulation, review, or lint did not catch it
At the pre-fix commit, Corundum's cocotb tests had no full-memory or X-state
check. Both used this queue range
\cite{corundum-parent,queue-manager-test,cpl-queue-manager-test}:

\begin{lstlisting}[language=Python]
QUEUE_INDEX_WIDTH = 8

for q in range(4):
    configure_queue(q)

q = random.randrange(4)
exercise_queue(q)
\end{lstlisting}

An 8-bit index provides 256 entries, but the tests used only queues 0 through
3. The faulty loop initialized entries 0 through 15, so every tested queue
appeared correct.

% - Was condition rare? property dont cover it? tool silent by construction?
The loop was legal syntax, so lint could not infer the missing addresses. The
fix added no full-range test or X-state assertion~\cite{corundum-fix}.
% -----------------------



% -----------------------
% 4. Class
\subsection{Classification}

% - Name of the class
This bug can be classified as a \textbf{Reset and initialization} bug.

% - Second instance of the same class (can be from section 4)
Rocket Chip provides a second independent instance of this class
\cite{rocket-init}. Its SRAMs were not initialized, which caused X propagation
in simulation.
% -----------------------



% -----------------------
% 5. Detection
\subsection{Detection}

% - SVA property/lint to catch this class
More generally, reset and initialization bugs can be detected by checking that
all relevant state reaches a valid value after reset. Lint can flag state that
is never reset, while simulation and formal properties can verify the expected
post-reset values.

\begin{lstlisting}
reset_complete |-> all_required_state_is_valid
\end{lstlisting}

For this reproducer, the concrete invariant is that every
queue-memory entry is known and zero.

The simulation implementation \cite{bug1-simulation-check}:

\begin{lstlisting}
#1;
for (int q = 0; q < QUEUE_COUNT; q++)
  assert (!$isunknown(queue_ram[q * 32 +: 32]) &&
          queue_ram[q * 32 +: 32] == '0);
\end{lstlisting}

The corresponding formal property is implemented in
\cite{bug1-formal-check}:

\begin{lstlisting}
property p_entry_is_initialized(logic [31:0] entry);
  @(posedge f_clk)
  !$isunknown(entry) && entry == '0;
endproperty

for (genvar q = 0; q < QUEUE_COUNT; q++) begin
  assert property (
    p_entry_is_initialized(f_queue_ram[q * 32 +: 32]));
end
\end{lstlisting}

% - What the above detection mechanism cost (false positives, runtime, required assumptions,
%   or design constraints)
Simulation is cheap for this small memory, and its runtime grows linearly with
the number of entries. Formal checking is exhaustive, but its state space
grows with the memory. Lint may also report state that is deliberately left
unreset, so its warnings need to be reviewed or waived.

Anvil's checker generally cannot detect this type of bug, and it accepts both
versions. Its generated RTL initializes the full register, masking this
particular bug~\cite{bug1-anvil-buggy,bug1-anvil-fixed,
bug1-anvil-results}.

% ----------------------------------------------------------------






















% ----------------------------------------------------------------
% Bug2
\bugsection{2}{Stale-data exposure from an empty FIFO}

% -----------------------


% -----------------------

% 1. Instance
\subsection{Instance}

% - Design
This bug comes from OpenTitan~\cite{opentitan}, an open-source silicon
root-of-trust project. It is a bit different from the other RTL bugs in the
report.

The FIFO buffers incoming UART words and exposes the previously received word
when it is empty. Since this is visible to software, an empty FIFO must not
reveal a byte previously received by another isolated application. The issue
therefore identifies this as a \textbf{security} concern.

Issue~\#2315~\cite{fifo-issue} tracked the report. PR~\#2420
documents and corrects the stale-data exposure~\cite{fifo-fix}, with the
change merged in commit~\cite{fifo-fix-commit}.
% -----------------------


% -----------------------

% 2. Mechanism
\subsection{Mechanism}

% - Cause
The FIFO reset its pointers but not its storage. Its output remained connected
to the stored word even when the FIFO was empty.

\begin{lstlisting}
// Buggy code
if (Pass)
  assign rdata = (fifo_empty && wvalid) ? wdata : storage_rdata;
else
  assign rdata = storage_rdata;
\end{lstlisting}

Therefore, removing the last byte cleared the valid state but left that byte
visible to the software. The fix masks the output with zero when the FIFO is empty.

\begin{lstlisting}
// Fixed code
assign rdata = empty ? '0 : rdata_int;
\end{lstlisting}

% - Minimal Repro
The buggy and fixed SystemVerilog variants are provided separately
\cite{bug2-sv-buggy,bug2-sv-fixed}, together with their Anvil versions
\cite{bug2-anvil-buggy,bug2-anvil-fixed}.

% -----------------------


% -----------------------

% 3. Why it survived
\subsection{Why It Survived}

% - Why simulation, review, or lint did not catch it
OpenTitan checked the FIFO output only while \texttt{rvalid} was high
\cite{fifo-parent,fifo-issue,fifo-fix}:

\begin{lstlisting}
assert property (rvalid |-> !$isunknown(rdata));
\end{lstlisting}

When the FIFO became empty, this assertion stopped checking \texttt{rdata}.
The UART tests also read data only from a non-empty FIFO
\cite{uart-testplan,uart-scoreboard}.

% - Was condition rare? property dont cover it? tool silent by construction?
The failing operation was therefore not tested. The fix added no test or
assertion~\cite{fifo-fix}.

% -----------------------


% -----------------------

% 4. Class
\subsection{Classification}

% - Name of the class
This bug can be classified as a \textbf{Buffering, FIFOs, and flow control} bug.

% - Second instance of the same class (can be from section 4)
The assignment already gives the OpenCores UART16550 as a second instance. It
used stale receive-FIFO data without first checking whether the FIFO was empty
\cite{uart16550-stale}.

% -----------------------


% -----------------------

% 5. Detection
\subsection{Detection}

% - SVA property/lint to catch this class
More generally, buffering and flow-control bugs can be detected by checking
that interface outputs and transfer decisions agree with the buffer's occupancy
state. Simulation and formal properties can check these relationships on every
cycle.

\begin{lstlisting}
transfer |-> valid && ready && occupancy_allows_transfer
\end{lstlisting}

For this reproducer, the concrete invariant is that the externally
visible FIFO output is zero whenever the FIFO is empty.

The simulation implementation \cite{bug2-simulation-check}:

\begin{lstlisting}
assert (empty && rdata == '0)
  else $fatal(1, "empty FIFO exposed %h", rdata);
\end{lstlisting}

The corresponding formal property is implemented in
\cite{bug2-formal-check}:

\begin{lstlisting}
always_ff @(posedge f_clk) begin
  if (!f_rst_n)
    f_past_valid <= 1'b0;
  else
    f_past_valid <= 1'b1;
end

property p_reset_before_past_valid;
  @(posedge f_clk)
  !f_past_valid |-> !f_rst_n;
endproperty: p_reset_before_past_valid

property p_empty_data_is_zero;
  @(posedge f_clk) disable iff (!f_rst_n)
  f_past_valid && f_empty |-> (f_rdata == '0);
endproperty: p_empty_data_is_zero

assume property (p_reset_before_past_valid);

a_empty_data_is_zero:
  assert property (p_empty_data_is_zero);
\end{lstlisting}

% - What the above detection mechanism cost (false positives, runtime, required assumptions,
%   or design constraints)
The check is cheap, but the fix also adds a data-width multiplexer.

Anvil's checker generally cannot detect this type of bug because the retained
storage is not modified while borrowed. It accepts both versions, so the
invariant must be checked separately
\cite{bug2-anvil-buggy,bug2-anvil-fixed,bug2-anvil-results}.

% ----------------------------------------------------------------

















% ----------------------------------------------------------------

% Bug3
\bugsection{3}{Unqualified asynchronous-FIFO payload crossing}

% -----------------------


% -----------------------

% 1. Instance
\subsection{Instance}

% - Design
This bug is from OpenTitan~\cite{opentitan}, in the asynchronous FIFO used
by its TileLink clock-domain-crossing (CDC) path.

The FIFO carries request and response payloads between the host and
device clocks. Its synchronized control shall ensure that destination-domain
logic observes the payload only after the source write is stable.

Issue~\#13834~\cite{cdc-issue} documented the unqualified asynchronous path.
The behavior was corrected in PR~\#13991~\cite{cdc-fix} and its fixing
commit~\cite{cdc-fix-commit}.

% -----------------------


% -----------------------

% 2. Mechanism
\subsection{Mechanism}

% - Cause
The asynchronous FIFO synchronized its valid indication into the read clock
domain, but exposed the associated multi-bit payload without gating it with
the synchronized valid signal.

\begin{lstlisting}
// Buggy code
assign rdata_o = rdata_int;
\end{lstlisting}

Hence, a source-domain write may toggle the destination-visible payload
during the synchronization delay. The fix drives the output net to zero while
the synchronized valid is low.

\begin{lstlisting}
// Fixed code
assign rdata_o = rvalid_o ? rdata_int : '0;
\end{lstlisting}

% - Minimal Repro
The buggy and fixed SystemVerilog variants are provided separately
\cite{bug3-sv-buggy,bug3-sv-fixed}, together with their Anvil versions
\cite{bug3-anvil-buggy,bug3-anvil-fixed}.

To keep the SystemVerilog reproducer minimal, both variants ignore writes after
the first. The Anvil versions retain repeated transfers because Anvil's channel
semantics express and check payload lifetimes with less code.

% -----------------------


% -----------------------

% 3. Why it survived
\subsection{Why It Survived}

% - Why simulation, review, or lint did not catch it
The pre-fix TL-UL testbench had no property for payload activity while
\texttt{rvalid\_o} was low~\cite{cdc-parent,tlul-dv}. A CDC tool found this unchecked
asynchronous path~\cite{cdc-issue,cdc-fix}.

% - Was condition rare? property dont cover it? tool silent by construction?
RTL simulation cannot model analog metastability. The fix changed two RTL
files and added no test or assertion~\cite{cdc-fix}.

% -----------------------


% -----------------------

% 4. Class
\subsection{Classification}

% - Name of the class
This bug can be classified as a \textbf{Clock-domain crossing and metastability}
bug.

% - Second instance of the same class (can be from section 4)
The assignment gives the historical OpenCores UART16550 receive path
as a second instance. It used the asynchronous receive input directly before
commit \texttt{780ac1c} added synchronizer flops~\cite{uart16550-cdc}.

% -----------------------


% -----------------------

% 5. Detection
\subsection{Detection}

% - SVA property/lint to catch this class
More generally, CDC bugs can be detected with structural analysis that
identifies unsynchronized paths and unsafe multi-bit crossings. Simulation and
formal properties can complement this analysis by checking how the destination
uses the crossing.

\begin{lstlisting}
reject unsynchronized control and unqualified payload crossings
\end{lstlisting}

For this reproducer, the concrete invariant is that the
destination payload remains zero until its synchronized valid signal is
asserted.

The simulation implementation \cite{bug3-simulation-check}:

\begin{lstlisting}
assert (rvalid || rdata == '0)
  else $fatal(1, "payload changed before valid");
\end{lstlisting}

The corresponding formal property is implemented in
\cite{bug3-formal-check}:

\begin{lstlisting}
always_ff @(posedge f_rd_clk) begin
  if (!f_rst_n)
    f_past_valid <= 1'b0;
  else
    f_past_valid <= 1'b1;
end

property p_reset_before_past_valid;
  @(posedge f_rd_clk)
  !f_past_valid |-> !f_rst_n;
endproperty: p_reset_before_past_valid

property p_invalid_payload_is_zero;
  @(posedge f_rd_clk) disable iff (!f_rst_n)
  f_past_valid && !f_rvalid |-> (f_rdata == '0);
endproperty: p_invalid_payload_is_zero

assume property (p_reset_before_past_valid);

a_invalid_payload_is_zero:
  assert property (p_invalid_payload_is_zero);
\end{lstlisting}

% - What the above detection mechanism cost (false positives, runtime, required assumptions,
%   or design constraints)
CDC analysis requires correct clock constraints and reviewed waivers. The
property is cheap but cannot model metastability, and the fix adds a multiplexer.

Anvil rejects the buggy channel model because it overwrites a register while an
outstanding send borrows it. The fixed model keeps the payload stable and is
accepted~\cite{bug3-anvil-buggy,bug3-anvil-fixed,bug3-anvil-results}.

% ----------------------------------------------------------------












% ----------------------------------------------------------------

% Bug4
\bugsection{4}{Delayed backpressure across a four-phase handshake}



% -----------------------

% 1. Instance
\subsection{Instance}

% - Design
This bug also comes from OpenTitan~\cite{opentitan}, in the entropy conditioner
and its SHA3/AES power-management handshake.

The conditioner collects a fixed message and requests that AES halt before
SHA3 processing to avoid a simultaneous power spike. Once this request begins,
upstream entropy must be backpressured until processing completes and the
four-phase request/acknowledge exchange returns to idle.

The handshaking delay was corrected in PR~\#15854~\cite{handshake-fix} and its
implementation commit~\cite{handshake-fix-commit}.

% -----------------------


% -----------------------

% 2. Mechanism
\subsection{Mechanism}

% - Cause
The conditioner used a registered ready mask to avoid a combinational loop,
but the mask remained open until the delayed SHA3 processing event.

\begin{lstlisting}
// Buggy code
assign sha3_msg_end = sha3_process;
mask_d = sha3_start ? 1'b1 : sha3_msg_end ? 1'b0 : mask_q;
\end{lstlisting}

Therefore, additional words could be accepted while the design waited for the
AES acknowledgement. The fix stops accepting new input when
\texttt{cs\_aes\_halt\_req} is asserted. It also waits for the acknowledgement
to return low before allowing the next operation.

\begin{lstlisting}
// Fixed code
assign sha3_msg_end = cs_aes_halt_req;
mask_d = sha3_start ? 1'b1 : sha3_msg_end ? 1'b0 : mask_q;
\end{lstlisting}

% - Minimal Repro
The buggy and fixed SystemVerilog variants are provided separately
\cite{bug4-sv-buggy,bug4-sv-fixed}, together with their Anvil versions
\cite{bug4-anvil-buggy,bug4-anvil-fixed}.

% -----------------------



% -----------------------

% 3. Why it survived
\subsection{Why It Survived}

% - Why simulation, review, or lint did not catch it
The UVM environment randomized AES handshake delays, but did not check that a
request immediately stopped message acceptance
\cite{handshake-parent,entropy-aes-agent,push-pull-agent,entropy-scoreboard}.

% - Was condition rare? property dont cover it? tool silent by construction?
The fix changed two RTL files and added no test or assertion
\cite{handshake-fix-commit}.

% -----------------------


% -----------------------

% 4. Class
\subsection{Classification}

% - Name of the class
This bug can be classified as a \textbf{Handshake and channel protocol} bug.

% - Second instance of the same class (can be from section 4)
The assignment already gives the Xilinx AXI-Lite demo slave as a second
instance. It accepted a transfer using conditions beyond the interface's
VALID-and-READY handshake~\cite{xilinx-axi-demo}.

% -----------------------


% -----------------------

% 5. Detection
\subsection{Detection}

% - SVA property/lint to catch this class
More generally, handshake bugs can be detected by checking when transfers are
accepted and by tracking each request through the required protocol sequence.
Simulation and formal properties can check both safety and eventual completion.

\begin{lstlisting}
accepted |-> valid && ready
request_started |-> required_sequence_holds_until_idle
\end{lstlisting}

For this reproducer, the concrete invariant is that the halt request immediately
blocks readiness.

The simulation implementation \cite{bug4-simulation-check}:

\begin{lstlisting}
assert (accepted == 2 && !msg_ready)
  else $fatal(1, "accepted input while waiting for ack");
\end{lstlisting}

The corresponding formal property is implemented in
\cite{bug4-formal-check}:

\begin{lstlisting}
always_ff @(posedge f_clk) begin
  if (!f_rst_n)
    f_past_valid <= 1'b0;
  else
    f_past_valid <= 1'b1;
end

property p_reset_before_past_valid;
  @(posedge f_clk)
  !f_past_valid |-> !f_rst_n;
endproperty: p_reset_before_past_valid

property p_reset_remains_deasserted;
  @(posedge f_clk)
  f_past_valid |-> f_rst_n;
endproperty: p_reset_remains_deasserted

property p_request_applies_backpressure;
  @(posedge f_clk) disable iff (!f_rst_n)
  f_past_valid && f_aes_req |-> !f_msg_ready;
endproperty: p_request_applies_backpressure

property p_ack_eventually_asserts;
  @(posedge f_clk) disable iff (!f_rst_n)
  f_past_valid && f_aes_req |->
    s_eventually f_aes_ack;
endproperty: p_ack_eventually_asserts

property p_ack_eventually_deasserts;
  @(posedge f_clk) disable iff (!f_rst_n)
  f_past_valid && f_aes_ack |->
    s_eventually !f_aes_ack;
endproperty: p_ack_eventually_deasserts

property p_handshake_eventually_returns_idle;
  @(posedge f_clk) disable iff (!f_rst_n)
  f_past_valid && f_aes_req |->
    s_eventually (!f_aes_req && !f_aes_ack);
endproperty: p_handshake_eventually_returns_idle

assume property (p_reset_before_past_valid);
assume property (p_reset_remains_deasserted);
assume property (p_ack_eventually_asserts);
assume property (p_ack_eventually_deasserts);

assert property (p_request_applies_backpressure);
assert property (p_handshake_eventually_returns_idle);
\end{lstlisting}

% - What the above detection mechanism cost (false positives, runtime, required assumptions,
%   or design constraints)
Unlike the other formal checks in this report, the handshake-completion check
is a liveness property. It requires fairness assumptions for the environment:
ACK must eventually rise after a request and eventually fall. Without these
assumptions, the tool can produce a counterexample in which the environment
holds ACK forever. The formal setup is available here~\cite{bug4-sby-checks}.

Liveness checks can be more expensive than safety checks. Insufficient fairness
assumptions may prevent a proof. Assumptions that are too strong or unrealistic
may make the result meaningless.

Anvil rejects the buggy channel model because it overwrites a borrowed payload
while the acknowledgement is outstanding. The fixed blocking-handshake model
is accepted~\cite{bug4-anvil-buggy,bug4-anvil-fixed,bug4-anvil-results}.

% ----------------------------------------------------------------











% ----------------------------------------------------------------
% Bug5
\bugsection{5}{CSR instruction incorrectly allowed to dual-issue}

% -----------------------


% -----------------------

% 1. Instance
\subsection{Instance}

% - Design
This bug comes from CVA6~\cite{cva6}, an open-source RISC-V processor, in its
superscalar issue stage.

The issue stage can dispatch two ordinary instructions if they are proven to
be independent. However, CSR operations access privileged state and can
redirect, trap, or change the context, so they must be serialized.

The issue restriction was corrected in PR~\#2329~\cite{cva6-csr}, through the
CSR-hazard commit~\cite{cva6-csr-hazard-commit} and CSR-port
commit~\cite{cva6-csr-port-commit}.

% -----------------------


% -----------------------

% 2. Mechanism
\subsection{Mechanism}

% - Cause
The issue-stage exclusion vector blocked only selected functional units when
slot 0 contained a CSR instruction.

\begin{lstlisting}
// Buggy code
CSR: begin
  fus_busy[1].alu       = 1'b1;
  fus_busy[1].ctrl_flow = 1'b1;
  fus_busy[1].csr       = 1'b1;
end
\end{lstlisting}

Therefore, other instruction classes could issue with a CSR, and a CSR could
also reach the second lane. The fix blocks every functional unit from the
second lane when slot 0 contains a CSR and ensures that a CSR uses only lane 0.

\begin{lstlisting}
// Fixed code
fus_busy[1].csr = 1'b1;

CSR: begin
  fus_busy[1] = '1;
end
\end{lstlisting}

% - Minimal Repro
The buggy and fixed SystemVerilog variants are provided separately
\cite{bug5-sv-buggy,bug5-sv-fixed}, together with their Anvil versions
\cite{bug5-anvil-buggy,bug5-anvil-fixed}.

% -----------------------


% -----------------------

% 3. Why it survived
\subsection{Why It Survived}

% - Why simulation, review, or lint did not catch it
The PR does not state how the bug was found. A visible failure requires a CSR
to issue beside another instruction and affect its behavior.



% - Was condition rare? property dont cover it? tool silent by construction?
CVA6 had CSR tests and an RVFI/UVM scoreboard
\cite{cva6-csr-tests,cva6-verification-scoreboard}. However, its default CI
target disabled superscalar issue and could not exercise the bug
\cite{cva6-parent,cva6-ci,cva6-config}.

Both fixes changed only \texttt{issue\_read\_operands.sv} and added
no test or assertion~\cite{cva6-csr-hazard-commit,cva6-csr-port-commit}.

% -----------------------


% -----------------------

% 4. Class
\subsection{Classification}

% - Name of the class
This bug can be classified as a \textbf{Pipeline control and hazards} bug.

% - Second instance of the same class (can be from section 4)
XiangShan provides a second instance of this class. Its rename-stage bypass
used the wrong dependency tag, so \texttt{csrr vl} immediately following
\texttt{vsetvli} read stale state~\cite{xiangshan-vl-bug,xiangshan-vl-fix}.

% -----------------------


% -----------------------

% 5. Detection
\subsection{Detection}

% - SVA property/lint to catch this class
More generally, pipeline-control and hazard bugs can be detected by checking
each instruction selected for issue against the design's dependency and
serialization rules. Simulation and formal properties can exercise different
combinations of instruction classes.

\begin{lstlisting}
issued(i) |-> dependencies_clear(i) && serialization_allows(i)
\end{lstlisting}

For this reproducer, the concrete invariant is that a CSR instruction
never issues beside another instruction and never uses the second issue lane.

The simulation implementation \cite{bug5-simulation-check}:

\begin{lstlisting}
assert (!issue1)
  else $fatal(1, "instruction issued beside CSR");
\end{lstlisting}

The corresponding formal properties are implemented in
\cite{bug5-formal-check}:

\begin{lstlisting}
property p_csr_is_serialized;
  @(posedge f_clk)
  f_issue0 && f_fu0 == CSR |-> !f_issue1;
endproperty: p_csr_is_serialized

property p_csr_uses_port_zero;
  @(posedge f_clk)
  !(f_issue1 && f_fu1 == CSR);
endproperty: p_csr_uses_port_zero
\end{lstlisting}

% - What the above detection mechanism cost (false positives, runtime, required assumptions,
%   or design constraints)
This pessimistic serialization avoids the bug, but it may be too restrictive
and limit performance. However, CSRs are generally rare and uncommon in hot
loops. The check itself is inexpensive.

Anvil's checker generally cannot detect this type of bug because it requires
high-level system reasoning.
It accepts both versions, so CSR serialization must be
checked separately~\cite{bug5-anvil-buggy,bug5-anvil-fixed,
bug5-anvil-results}.

% ----------------------------------------------------------------










% ----------------------------------------------------------------
% Synthesis
\onecolumn
\section{Synthesis}
\titlespacing*{\subsection}{0pt}{1.0ex plus 0.3ex minus 0.2ex}{0.35ex}

\begingroup
\centering
\footnotesize
\setlength{\tabcolsep}{3pt}
\renewcommand{\arraystretch}{1.0}
\begin{tabularx}{\textwidth}{c p{0.20\textwidth} X p{0.20\textwidth} p{0.16\textwidth}}
\hline
\textbf{Bug} & \textbf{Class} & \textbf{Comments} &
\textbf{Best technique} & \textbf{Anvil ability} \\
\hline
1 & Reset and initialization &
The loop initializes only a part of a parameterized memory. &
Four-state simulation or a formal completeness property over every entry. &
Not detected. The generated RTL masks the difference. \\
\hline

2 & Buffering, FIFOs, and flow control &
An empty software-visible FIFO exposes its retained payload. &
Empty=>zero assertion. &
Not detected. \\
\hline

3 & Clock-domain crossing and metastability &
An asynchronous payload is visible before its synchronized valid signal. &
Structural CDC analysis plus an invalid=>zero assertion. &
Detected when modeled as a channel-lifetime violation. \\
\hline

4 & Handshake and channel protocol &
Delayed backpressure accepts extra input while acknowledgement is pending. &
Protocol assertions with continuous input and delayed ACK. &
Detected when modeled as a channel-lifetime violation. \\
\hline

5 & Pipeline control and hazards &
A CSR can issue beside another instruction or on the wrong port. &
Issue serialization assertions and formal instruction pairing. &
Not detected. \\
\hline
\end{tabularx}
\par
\endgroup

% ------------------


\subsection{Anvil Modeling Requirement}

For Bugs 3 and 4, Anvil rejects the buggy versions when the designs are
correctly expressed using its channel model
\cite{bug3-anvil-buggy,bug4-anvil-buggy}. The corrected versions pass the
compiler~\cite{bug3-anvil-fixed,bug4-anvil-fixed}. However, code that models
the channel incorrectly can still pass Anvil's checks while remaining
functionally wrong~\cite{bug3-anvil-semantic,bug4-anvil-semantic}.

% ------------------

\subsection{Simulation}

The simulations use the open-source SystemVerilog/Verilog simulator
Verilator 5.051, revision \texttt{7d3021c34}~\cite{verilator-version}.

Verilator is a two-state simulator. It is also cycle-based, although that is
not relevant here. 
Hence, uninitialized values are not preserved as \texttt{X}, so both Bug~1 versions pass.


% ------------------

\subsection{Model Checking}

The formal checks use the open-source SymbiYosys tool, version 0.67,
revision \texttt{fea6e46}~\cite{symbiyosys-version}. BMC finds
counterexamples in the buggy versions. The fixed versions use $k$-induction
for Bug~1 and PDR for Bugs~2--5. Bug~4 also checks liveness with
\texttt{suprove}.
% ------------------

\subsection{CDC}

For Bug~3, rtl-buddy-cdc reported the buggy payload as an unprotected bus
crossing. It did not report this issue for the fixed, valid-gated
payload~\cite{rtl-buddy-cdc,bug3-cdc-results}.

% ------------------

\subsection{Selected Projects}

The selected projects are Corundum~\cite{corundum},
OpenTitan~\cite{opentitan}, and CVA6~\cite{cva6}.

\begin{itemize}[nosep]
\item OpenTitan includes a UVM and formal-verification suite, and the selected bugs have
  clear issue-PR-commit links.
\item CVA6 includes processor, ISA, and UVM verification infrastructure.
\item Corundum uses cocotb/pytest regressions and does not include UVM or
  formal checks.
\end{itemize}

\subsection{Test Matrix}

\begingroup
\centering
\rotatebox{90}{\begin{minipage}{0.84\textheight}
\scriptsize
\setlength{\tabcolsep}{0.8pt}
\renewcommand{\arraystretch}{1.15}
\begin{tabularx}{\linewidth}{c *{13}{>{\raggedright\arraybackslash}X}}
\hline
\textbf{Bug} &
\shortstack{\textbf{Sim/SV}\\\textbf{buggy}\\\cite{matrix-sv-sim}} &
\shortstack{\textbf{Sim/SV}\\\textbf{fixed}\\\cite{matrix-sv-sim}} &
\shortstack{\textbf{Formal/SV}\\\textbf{buggy}\\\cite{matrix-sv-formal}} &
\shortstack{\textbf{Formal/SV}\\\textbf{fixed}\\\cite{matrix-sv-formal}} &
\shortstack{\textbf{Compile}\\\textbf{Anvil}\\\textbf{buggy}\\\cite{matrix-anvil-sources}} &
\shortstack{\textbf{Compile}\\\textbf{Anvil}\\\textbf{semantic}\\\cite{matrix-anvil-semantic-sources}} &
\shortstack{\textbf{Compile}\\\textbf{Anvil}\\\textbf{fixed}\\\cite{matrix-anvil-sources}} &
\shortstack{\textbf{Sim/Anvil}\\\textbf{buggy}\\\cite{matrix-anvil-sim}} &
\shortstack{\textbf{Sim/Anvil}\\\textbf{semantic}\\\cite{matrix-anvil-sim}} &
\shortstack{\textbf{Sim/Anvil}\\\textbf{fixed}\\\cite{matrix-anvil-sim}} &
\shortstack{\textbf{Formal}\\\textbf{Anvil}\\\textbf{buggy}\\\cite{matrix-anvil-formal}} &
\shortstack{\textbf{Formal}\\\textbf{Anvil}\\\textbf{semantic}\\\cite{matrix-anvil-formal}} &
\shortstack{\textbf{Formal}\\\textbf{Anvil}\\\textbf{fixed}\\\cite{matrix-anvil-formal}} \\
\hline
1 & PASS & PASS &
Safety; BMC; Z3; d2; FAIL/SAT &
Safety; BMC+$k$-ind.; Z3; d2; PASS &
PASS & --- & PASS &
PASS (masked) & --- & PASS &
Safety; BMC; Z3; d6; PASS (masked) & --- &
Safety; BMC+$k$-ind.; Z3; d6; PASS \\
\hline

2 & FAIL & PASS &
Safety; BMC; Z3; d8; FAIL/SAT &
Safety; PDR; ABC; PASS &
PASS & --- & PASS &
FAIL & --- & PASS &
Safety; BMC; Z3; d12; FAIL/SAT & --- &
Safety; PDR; ABC; PASS \\
\hline

3 & FAIL & PASS &
Safety; BMC; Z3; d8; FAIL/SAT &
Safety; PDR; ABC; PASS &
REJECT & PASS & PASS &
Not run & FAIL & PASS &
Not run & Safety; BMC; Z3; d12; FAIL/SAT &
Safety; PDR; ABC; PASS \\
\hline

4-S & FAIL & PASS &
Safety; BMC; Z3; d10; FAIL/SAT &
Safety; PDR; ABC; PASS &
REJECT & PASS & PASS &
Not run & FAIL & PASS &
Not run & Safety; BMC; Z3; d16; FAIL/SAT &
Safety; PDR; ABC; PASS \\
\hline

4-L & --- & --- &
Live; \texttt{suprove}; PASS &
Live; \texttt{suprove}; PASS &
--- & --- & --- & --- & --- & --- &
--- & --- & Live; \texttt{suprove}; PASS \\
\hline

5 & FAIL & PASS &
Safety; BMC; Z3; d2; FAIL/SAT &
Safety; PDR; ABC; PASS &
PASS & --- & PASS &
FAIL & --- & PASS &
Safety; BMC; Z3; d12; FAIL/SAT & --- &
Safety; PDR; ABC; PASS \\
\hline
\end{tabularx}
\end{minipage}}
\par
\endgroup
% ----------------------------------------------------------------


% References
\titlespacing*{\subsection}{0pt}{1.8ex plus 0.5ex minus 0.2ex}{0.7ex}
\twocolumn
\begin{thebibliography}{99}

\bibitem{matrix-sv-sim}
SystemVerilog simulation environments:
\href{run:../src/bug1}{\path{src/bug1}},
\href{run:../src/bug2}{\path{src/bug2}},
\href{run:../src/bug3}{\path{src/bug3}},
\href{run:../src/bug4}{\path{src/bug4}},
\href{run:../src/bug5}{\path{src/bug5}}.

\bibitem{matrix-sv-formal}
SystemVerilog formal environments:
\href{run:../src/bug1/formal}{\path{src/bug1/formal}},
\href{run:../src/bug2/formal}{\path{src/bug2/formal}},
\href{run:../src/bug3/formal}{\path{src/bug3/formal}},
\href{run:../src/bug4/formal}{\path{src/bug4/formal}},
\href{run:../src/bug5/formal}{\path{src/bug5/formal}}.

\bibitem{matrix-anvil-sources}
Buggy and fixed Anvil sources:
\href{run:../src/bug1/rtl}{\path{src/bug1/rtl}},
\href{run:../src/bug2/rtl}{\path{src/bug2/rtl}},
\href{run:../src/bug3/rtl}{\path{src/bug3/rtl}},
\href{run:../src/bug4/rtl}{\path{src/bug4/rtl}},
\href{run:../src/bug5/rtl}{\path{src/bug5/rtl}}.

\bibitem{matrix-anvil-semantic-sources}
Anvil semantic-bug sources:
\href{run:../src/bug3/rtl/semantic_bug.anvil}{\path{src/bug3/rtl/semantic_bug.anvil}},
\href{run:../src/bug4/rtl/semantic_bug.anvil}{\path{src/bug4/rtl/semantic_bug.anvil}}.

\bibitem{matrix-anvil-sim}
Generated-Anvil simulation testbenches:
\href{run:../src/bug1/tb/anvil_tb.sv}{\path{src/bug1/tb/anvil_tb.sv}},
\href{run:../src/bug2/tb/anvil_tb.sv}{\path{src/bug2/tb/anvil_tb.sv}},
\href{run:../src/bug3/tb/anvil_tb.sv}{\path{src/bug3/tb/anvil_tb.sv}},
\href{run:../src/bug4/tb/anvil_tb.sv}{\path{src/bug4/tb/anvil_tb.sv}}, and
\href{run:../src/bug5/tb/anvil_tb.sv}{\path{src/bug5/tb/anvil_tb.sv}}.

\bibitem{matrix-anvil-formal}
Generated-Anvil formal configurations, properties, and logs:
\href{run:../src/bug1/formal}{\path{src/bug1/formal}},
\href{run:../src/bug2/formal}{\path{src/bug2/formal}},
\href{run:../src/bug3/formal}{\path{src/bug3/formal}},
\href{run:../src/bug4/formal}{\path{src/bug4/formal}},
\href{run:../src/bug5/formal}{\path{src/bug5/formal}}.

\bibitem{verilator-version}
\href{https://github.com/verilator/verilator/commit/7d3021c34d0f26a88e8b1107c50ceb2d896689b3}{Verilator revision
\path{7d3021c34d0f26a88e8b1107c50ceb2d896689b3}}.

\bibitem{symbiyosys-version}
\href{https://github.com/YosysHQ/sby/commit/fea6e467d067b3ea84b6b5ac08cd48beb59f0d42}{SymbiYosys version 0.67, revision
\path{fea6e467d067b3ea84b6b5ac08cd48beb59f0d42}}.

\bibitem{rtl-buddy-cdc}
\href{https://github.com/rtl-buddy/rtl-buddy-cdc/commit/df996d55af0e8b824183d03d531426b8601af429}{rtl-buddy-cdc version 0.5.0, revision
\path{df996d55af0e8b824183d03d531426b8601af429}}.

\bibitem{corundum}
\href{https://github.com/corundum/corundum}{Corundum source repository}.

\bibitem{queue-manager-src}
\href{https://github.com/corundum/corundum/blob/0560f98e799d741d62522e61bf23321fc3f2880b/fpga/common/rtl/queue_manager.v}{Corundum \texttt{queue\_manager} source}.

\bibitem{cpl-queue-manager-src}
\href{https://github.com/corundum/corundum/blob/0560f98e799d741d62522e61bf23321fc3f2880b/fpga/common/rtl/cpl_queue_manager.v}{Corundum \texttt{cpl\_queue\_manager} source}.

\bibitem{tx-scheduler-src}
\href{https://github.com/corundum/corundum/blob/0560f98e799d741d62522e61bf23321fc3f2880b/fpga/common/rtl/tx_scheduler_rr.v}{Corundum \texttt{tx\_scheduler\_rr} source}.

\bibitem{queue-manager-test}
\href{https://github.com/corundum/corundum/blob/0f82e0c5f3b9fa9a15e99d6e43b08b3b4dff7909/fpga/common/tb/queue_manager/test_queue_manager.py}{Corundum pre-fix \texttt{queue\_manager} regression test}.

\bibitem{cpl-queue-manager-test}
\href{https://github.com/corundum/corundum/blob/0f82e0c5f3b9fa9a15e99d6e43b08b3b4dff7909/fpga/common/tb/cpl_queue_manager/test_cpl_queue_manager.py}{Corundum pre-fix \texttt{cpl\_queue\_manager} regression test}.

\bibitem{corundum-intro}
\href{https://github.com/corundum/corundum/commit/0560f98e799d741d62522e61bf23321fc3f2880b}{Corundum: Support more than 4k queues}.

\bibitem{corundum-fix}
\href{https://github.com/corundum/corundum/pull/96}{Corundum PR~\#96: Correct queue-memory initialization}.

\bibitem{corundum-fix-commit}
\href{https://github.com/corundum/corundum/commit/7bc1ecc711f3d4fce55290e0ed6d94a549cc2528}{Corundum fixing commit \texttt{7bc1ecc}}.

\bibitem{corundum-parent}
\href{https://github.com/corundum/corundum/commit/0f82e0c5f3b9fa9a15e99d6e43b08b3b4dff7909}{Corundum pre-fix parent commit \texttt{0f82e0c}}.

\bibitem{quartus}
\href{https://www.intel.com/content/www/us/en/docs/programmable/683472/24-3/design-suite-overview.html}{Intel Quartus Prime FPGA design software}.

\bibitem{rocket-init}
\href{https://github.com/chipsalliance/rocket-chip/commit/1cb2d1d7b7191aff8f8004b9b64a5d908c96fd6b}{Rocket Chip: Initialize all SRAMs to avoid X propagation}.

\bibitem{opentitan}
\href{https://github.com/lowRISC/opentitan}{OpenTitan source repository}.

\bibitem{fifo-issue}
\href{https://github.com/lowRISC/opentitan/issues/2315}{OpenTitan issue~\#2315: Stale data from empty UART RX FIFO}.

\bibitem{fifo-fix}
\href{https://github.com/lowRISC/opentitan/pull/2420}{OpenTitan PR~\#2420: Fix empty-FIFO stale data}.

\bibitem{fifo-fix-commit}
\href{https://github.com/lowRISC/opentitan/commit/9ac2b15f162cc4bf5e6e4226ae659c9d4225e1d0}{OpenTitan merged commit \texttt{9ac2b15}}.

\bibitem{fifo-parent}
\href{https://github.com/lowRISC/opentitan/commit/0408a4e7564c0fde1bf7dca87e8c8e1bce587c75}{OpenTitan stale-FIFO pre-fix parent commit \texttt{0408a4e}}.

\bibitem{uart-testplan}
\href{https://github.com/lowRISC/opentitan/blob/0408a4e7564c0fde1bf7dca87e8c8e1bce587c75/hw/ip/uart/data/uart_testplan.hjson}{OpenTitan pre-fix UART test plan}.

\bibitem{uart-scoreboard}
\href{https://github.com/lowRISC/opentitan/blob/0408a4e7564c0fde1bf7dca87e8c8e1bce587c75/hw/ip/uart/dv/env/uart_scoreboard.sv}{OpenTitan pre-fix UART scoreboard}.

\bibitem{cdc-issue}
\href{https://github.com/lowRISC/opentitan/issues/13834}{OpenTitan issue~\#13834: W\_MASYNC in spi\_device}.

\bibitem{cdc-fix}
\href{https://github.com/lowRISC/opentitan/pull/13991}{OpenTitan PR~\#13991: Qualify asynchronous-FIFO output}.

\bibitem{cdc-fix-commit}
\href{https://github.com/lowRISC/opentitan/commit/e0ba6cdb527c02b7e85dff048a8f168a6e166999}{OpenTitan CDC-fix commit \texttt{e0ba6cd}}.

\bibitem{cdc-parent}
\href{https://github.com/lowRISC/opentitan/commit/592eaf648de3fdc0803761f9fb31390012587573}{OpenTitan CDC pre-fix parent commit \texttt{592eaf6}}.

\bibitem{tlul-dv}
\href{https://github.com/lowRISC/opentitan/blob/592eaf648de3fdc0803761f9fb31390012587573/hw/ip/tlul/doc/dv/index.md}{OpenTitan pre-fix TL-UL verification environment}.

\bibitem{handshake-fix}
\href{https://github.com/lowRISC/opentitan/pull/15854}{OpenTitan PR~\#15854: Fix handshaking delay}.

\bibitem{handshake-fix-commit}
\href{https://github.com/lowRISC/opentitan/commit/cd1b6131548f8bdb86e2f10904d0b8551db31d50}{OpenTitan handshake-fix commit \texttt{cd1b613}}.

\bibitem{handshake-parent}
\href{https://github.com/lowRISC/opentitan/commit/78944e3d4b0af31133ccca88395de4a74d0c6b8b}{OpenTitan handshake pre-fix parent commit \texttt{78944e3}}.

\bibitem{entropy-aes-agent}
\href{https://github.com/lowRISC/opentitan/blob/78944e3d4b0af31133ccca88395de4a74d0c6b8b/hw/ip/entropy_src/dv/env/entropy_src_env.sv}{OpenTitan pre-fix entropy-source AES handshake agent}.

\bibitem{push-pull-agent}
\href{https://github.com/lowRISC/opentitan/blob/78944e3d4b0af31133ccca88395de4a74d0c6b8b/hw/dv/sv/push_pull_agent/push_pull_agent_cfg.sv}{OpenTitan pre-fix randomized handshake-delay configuration}.

\bibitem{entropy-scoreboard}
\href{https://github.com/lowRISC/opentitan/blob/78944e3d4b0af31133ccca88395de4a74d0c6b8b/hw/ip/entropy_src/dv/env/entropy_src_scoreboard.sv}{OpenTitan pre-fix entropy-source scoreboard}.

\bibitem{cva6}
\href{https://github.com/openhwfoundation/cva6}{CVA6 source repository}.

\bibitem{cva6-csr}
\href{https://github.com/openhwfoundation/cva6/pull/2329}{CVA6 PR~\#2329: Do not issue CSR with another instruction}.

\bibitem{cva6-csr-hazard-commit}
\href{https://github.com/openhwfoundation/cva6/commit/92e418a507134ef4cb95925feba4a72ff7f845cb}{CVA6 CSR-hazard commit \texttt{92e418a}}.

\bibitem{cva6-csr-port-commit}
\href{https://github.com/openhwfoundation/cva6/commit/827daa09f413950e28e4f4cb84843e9d8e41a4c5}{CVA6 CSR-port commit \texttt{827daa0}}.

\bibitem{cva6-parent}
\href{https://github.com/openhwfoundation/cva6/commit/67dba2cad376de0190bb5564f1e229de73e86ef1}{CVA6 pre-fix parent commit \texttt{67dba2c}}.

\bibitem{cva6-ci}
\href{https://github.com/openhwfoundation/cva6/blob/67dba2cad376de0190bb5564f1e229de73e86ef1/.gitlab-ci.yml}{CVA6 pre-fix CI configuration}.

\bibitem{cva6-config}
\href{https://github.com/openhwfoundation/cva6/blob/67dba2cad376de0190bb5564f1e229de73e86ef1/core/include/cv32a65x_config_pkg.sv}{CVA6 pre-fix \texttt{cv32a65x} configuration}.

\bibitem{cva6-csr-tests}
\href{https://github.com/openhwfoundation/cva6/blob/67dba2cad376de0190bb5564f1e229de73e86ef1/verif/regress/dv-riscv-csr-access-test.sh}{CVA6 pre-fix CSR regression}.

\bibitem{cva6-verification-scoreboard}
\href{https://github.com/openhwfoundation/cva6/blob/67dba2cad376de0190bb5564f1e229de73e86ef1/verif/env/uvme/uvme_cva6_sb.sv}{CVA6 pre-fix RVFI/UVM scoreboard}.

\bibitem{uart16550-stale}
\href{https://github.com/freecores/uart16550/blob/master/rtl/verilog/uart_regs.v}{OpenCores UART16550 receive-FIFO RTL}.

\bibitem{uart16550-cdc}
\href{https://github.com/freecores/uart16550/commit/780ac1c1252506fe719f306f6c30364a8b1e01eb}{OpenCores UART16550 commit \texttt{780ac1c}: add receive-input synchronizer flops}.

\bibitem{xiangshan-vl-bug}
\href{https://github.com/OpenXiangShan/XiangShan/issues/5739}{XiangShan issue~\#5739: \texttt{csrr vl} reads stale zero immediately after \texttt{vsetvli}}.

\bibitem{xiangshan-vl-fix}
\href{https://github.com/OpenXiangShan/XiangShan/pull/5743}{XiangShan PR~\#5743: fix vector-length dependency bypass in rename}.

\bibitem{intel-errata}
\href{https://www.intel.com/content/www/us/en/content-details/334408/}{Intel processor specification update: SKL150/KBL095}.

\bibitem{xilinx-axi-demo}
\href{https://zipcpu.com/formal/2019/05/13/axifull.html}{D. Gisselquist,
``Examining Xilinx's AXI demonstration core.''}

\bibitem{anvil-paper}
J. Z. Yu, A. R. Jha, U. Mathur, T. E. Carlson, and P. Saxena,
\href{https://arxiv.org/abs/2503.19447}{``Anvil: A General-Purpose Timing-Safe Hardware Description Language,''}
ASPLOS, 2026.

\bibitem{anvil-github}
\href{https://github.com/AnvilHDL/anvil}{Anvil compiler source repository}.

% 2. Repro/
\bibitem{bug1-sv-buggy}
\href{run:../src/bug1/rtl/buggy.sv}{\path{src/bug1/rtl/buggy.sv}}.

\bibitem{bug1-anvil-buggy}
\href{run:../src/bug1/rtl/buggy.anvil}{\path{src/bug1/rtl/buggy.anvil}}.

\bibitem{bug1-sv-fixed}
\href{run:../src/bug1/rtl/fixed.sv}{\path{src/bug1/rtl/fixed.sv}}.

\bibitem{bug1-anvil-fixed}
\href{run:../src/bug1/rtl/fixed.anvil}{\path{src/bug1/rtl/fixed.anvil}}.

\bibitem{bug1-simulation-check}
\href{run:../src/bug1/tb/tb.sv}{\path{src/bug1/tb/tb.sv}}.

\bibitem{bug1-formal-check}
\href{run:../src/bug1/formal/properties.sv}{\path{src/bug1/formal/properties.sv}}.

\bibitem{bug1-anvil-results}
\href{run:../src/bug1/formal/results/logs/anvil_compile_buggy.log}{\path{src/bug1/formal/results/logs/anvil_compile_buggy.log}};
\href{run:../src/bug1/formal/results/logs/anvil_compile_fixed.log}{\path{src/bug1/formal/results/logs/anvil_compile_fixed.log}};
\href{run:../src/bug1/formal/results/logs/anvil_buggy.log}{\path{src/bug1/formal/results/logs/anvil_buggy.log}};
\href{run:../src/bug1/formal/results/logs/anvil_fixed.log}{\path{src/bug1/formal/results/logs/anvil_fixed.log}}.

\bibitem{bug2-sv-buggy}
\href{run:../src/bug2/rtl/buggy.sv}{\path{src/bug2/rtl/buggy.sv}}.

\bibitem{bug2-sv-fixed}
\href{run:../src/bug2/rtl/fixed.sv}{\path{src/bug2/rtl/fixed.sv}}.

\bibitem{bug2-anvil-buggy}
\href{run:../src/bug2/rtl/buggy.anvil}{\path{src/bug2/rtl/buggy.anvil}}.

\bibitem{bug2-anvil-fixed}
\href{run:../src/bug2/rtl/fixed.anvil}{\path{src/bug2/rtl/fixed.anvil}}.

\bibitem{bug2-simulation-check}
\href{run:../src/bug2/tb/tb.sv}{\path{src/bug2/tb/tb.sv}}.

\bibitem{bug2-formal-check}
\href{run:../src/bug2/formal/properties.sv}{\path{src/bug2/formal/properties.sv}}.

\bibitem{bug2-anvil-results}
\href{run:../src/bug2/formal/results/logs/anvil_compile_buggy.log}{\path{src/bug2/formal/results/logs/anvil_compile_buggy.log}};
\href{run:../src/bug2/formal/results/logs/anvil_compile_fixed.log}{\path{src/bug2/formal/results/logs/anvil_compile_fixed.log}};
\href{run:../src/bug2/formal/results/logs/anvil_buggy.log}{\path{src/bug2/formal/results/logs/anvil_buggy.log}};
\href{run:../src/bug2/formal/results/logs/anvil_fixed.log}{\path{src/bug2/formal/results/logs/anvil_fixed.log}}.

\bibitem{bug3-sv-buggy}
\href{run:../src/bug3/rtl/buggy.sv}{\path{src/bug3/rtl/buggy.sv}}.

\bibitem{bug3-sv-fixed}
\href{run:../src/bug3/rtl/fixed.sv}{\path{src/bug3/rtl/fixed.sv}}.

\bibitem{bug3-anvil-buggy}
\href{run:../src/bug3/rtl/buggy.anvil}{\path{src/bug3/rtl/buggy.anvil}}.

\bibitem{bug3-anvil-fixed}
\href{run:../src/bug3/rtl/fixed.anvil}{\path{src/bug3/rtl/fixed.anvil}}.

\bibitem{bug3-anvil-semantic}
\href{run:../src/bug3/rtl/semantic_bug.anvil}{\path{src/bug3/rtl/semantic_bug.anvil}}.

\bibitem{bug3-simulation-check}
\href{run:../src/bug3/tb/tb.sv}{\path{src/bug3/tb/tb.sv}}.

\bibitem{bug3-formal-check}
\href{run:../src/bug3/formal/properties.sv}{\path{src/bug3/formal/properties.sv}}.

\bibitem{bug3-anvil-results}
\href{run:../src/bug3/formal/results/logs/anvil_compile_buggy.log}{\path{src/bug3/formal/results/logs/anvil_compile_buggy.log}};
\href{run:../src/bug3/formal/results/logs/anvil_compile_semantic_bug.log}{\path{src/bug3/formal/results/logs/anvil_compile_semantic_bug.log}};
\href{run:../src/bug3/formal/results/logs/anvil_compile_fixed.log}{\path{src/bug3/formal/results/logs/anvil_compile_fixed.log}};
\href{run:../src/bug3/formal/results/logs/anvil_semantic_bug.log}{\path{src/bug3/formal/results/logs/anvil_semantic_bug.log}};
\href{run:../src/bug3/formal/results/logs/anvil_fixed.log}{\path{src/bug3/formal/results/logs/anvil_fixed.log}}.

\bibitem{bug3-cdc-results}
\href{run:../src/bug3/cdc/results/buggy.log}{\path{src/bug3/cdc/results/buggy.log}};
\href{run:../src/bug3/cdc/results/fixed.log}{\path{src/bug3/cdc/results/fixed.log}}.

\bibitem{bug4-sv-buggy}
\href{run:../src/bug4/rtl/buggy.sv}{\path{src/bug4/rtl/buggy.sv}}.

\bibitem{bug4-sv-fixed}
\href{run:../src/bug4/rtl/fixed.sv}{\path{src/bug4/rtl/fixed.sv}}.

\bibitem{bug4-anvil-buggy}
\href{run:../src/bug4/rtl/buggy.anvil}{\path{src/bug4/rtl/buggy.anvil}}.

\bibitem{bug4-anvil-fixed}
\href{run:../src/bug4/rtl/fixed.anvil}{\path{src/bug4/rtl/fixed.anvil}}.

\bibitem{bug4-anvil-semantic}
\href{run:../src/bug4/rtl/semantic_bug.anvil}{\path{src/bug4/rtl/semantic_bug.anvil}}.

\bibitem{bug4-simulation-check}
\href{run:../src/bug4/tb/tb.sv}{\path{src/bug4/tb/tb.sv}}.

\bibitem{bug4-formal-check}
\href{run:../src/bug4/formal/properties.sv}{\path{src/bug4/formal/properties.sv}}.

\bibitem{bug4-sby-checks}
\href{run:../src/bug4/formal/sby_properties.sv}{\path{src/bug4/formal/sby_properties.sv}};
\href{run:../src/bug4/formal.sby}{\path{src/bug4/formal.sby}};
\href{run:../src/bug4/formal_live.sby}{\path{src/bug4/formal_live.sby}}.

\bibitem{bug4-anvil-results}
\href{run:../src/bug4/formal/results/logs/anvil_compile_buggy.log}{\path{src/bug4/formal/results/logs/anvil_compile_buggy.log}};
\href{run:../src/bug4/formal/results/logs/anvil_compile_semantic_bug.log}{\path{src/bug4/formal/results/logs/anvil_compile_semantic_bug.log}};
\href{run:../src/bug4/formal/results/logs/anvil_compile_fixed.log}{\path{src/bug4/formal/results/logs/anvil_compile_fixed.log}};
\href{run:../src/bug4/formal/results/logs/anvil_semantic_bug.log}{\path{src/bug4/formal/results/logs/anvil_semantic_bug.log}};
\href{run:../src/bug4/formal/results/logs/anvil_fixed.log}{\path{src/bug4/formal/results/logs/anvil_fixed.log}}.

\bibitem{bug5-sv-buggy}
\href{run:../src/bug5/rtl/buggy.sv}{\path{src/bug5/rtl/buggy.sv}}.

\bibitem{bug5-sv-fixed}
\href{run:../src/bug5/rtl/fixed.sv}{\path{src/bug5/rtl/fixed.sv}}.

\bibitem{bug5-anvil-buggy}
\href{run:../src/bug5/rtl/buggy.anvil}{\path{src/bug5/rtl/buggy.anvil}}.

\bibitem{bug5-anvil-fixed}
\href{run:../src/bug5/rtl/fixed.anvil}{\path{src/bug5/rtl/fixed.anvil}}.

\bibitem{bug5-simulation-check}
\href{run:../src/bug5/tb/tb.sv}{\path{src/bug5/tb/tb.sv}}.

\bibitem{bug5-formal-check}
\href{run:../src/bug5/formal/properties.sv}{\path{src/bug5/formal/properties.sv}}.

\bibitem{bug5-anvil-results}
\href{run:../src/bug5/formal/results/logs/anvil_compile_buggy.log}{\path{src/bug5/formal/results/logs/anvil_compile_buggy.log}};
\href{run:../src/bug5/formal/results/logs/anvil_compile_fixed.log}{\path{src/bug5/formal/results/logs/anvil_compile_fixed.log}};
\href{run:../src/bug5/formal/results/logs/anvil_buggy.log}{\path{src/bug5/formal/results/logs/anvil_buggy.log}};
\href{run:../src/bug5/formal/results/logs/anvil_fixed.log}{\path{src/bug5/formal/results/logs/anvil_fixed.log}}.
\end{thebibliography}

\section*{Revision History}

\begin{itemize}
  \item 1: Replaced the second example for Bug~5, Intel's partial-register
  erratum~\cite{intel-errata}, with the more relevant XiangShan dependency
  bug~\cite{xiangshan-vl-bug,xiangshan-vl-fix}.
  \item 2: Generalized the detection discussions for Bugs~1-5.
  \item 3: Clarified Bug~3's reduction scope.
\end{itemize}

\end{document}
